%======================================================================%
\section{Proton Decay from the $\E{6}$ Breaking Chain}
\label{sec:proton-decay}
%======================================================================%

Prediction~\ref{pred:proton-decay} records a conservative lower bound on
$\tau(p\to e^+\pi^0)$ from the D3 breaking chain. This section closes the
contact by deriving the \emph{saturation} lifetime: the rate at which the
dimension-six $\mathrm{SU}(5)$ operator induced by $\mathbf{45}_H$ is fully
expressed once $M_{\mathrm{GUT}}$, $\alpha_{\mathrm{GUT}}$, and observer
overlap data are fixed by Theorem~\ref{thm:breaking-complete},
Proposition~\ref{prop:gut-scale}, and Theorem~\ref{thm:duality}. No
additional parameters enter beyond the capacity weights already present in
the Yukawa programme (Section~\ref{sec:yukawa}).

%----------------------------------------------------------------------%
\subsection{Dimension-six operator from staged breaking}
\label{sec:proton-decay:operator}
%----------------------------------------------------------------------%

At the $\iota=3$ stage of Theorem~\ref{thm:breaking-complete}, gradient flow
along $\Phi_{\mathbf{45}}\in\mathbf{45}_H$ reaches the intermediate vacuum
$\Spin(10)\to\mathrm{SU}(5)\times U(1)_X$
(Theorem~\ref{thm:breaking-so10su5}). The adjoint VEV generates heavy
$X_\mu,Y_\mu$ gauge bosons with mass
\begin{equation}\label{eq:proton-XY-mass}
  M_X \;\approx\; M_Y \;\approx\; M_{\mathrm{GUT}},
\end{equation}
with $M_{\mathrm{GUT}}$ identified by Proposition~\ref{prop:gut-scale}. Below
$M_{\mathrm{GUT}}$, the leading baryon-number-violating effective Lagrangian is
the dimension-six $\mathrm{SU}(5)$ operator
\begin{equation}\label{eq:proton-dim6}
  \mathcal{L}_{B=1}^{(6)}
  \;=\;
  \frac{g_{\mathrm{eff}}^{2}}{M_{\mathrm{GUT}}^{2}}\,
  \bigl(\overline{q}\,\overline{q}\,\overline{q}\,\ell\bigr)
  \;+\;\mathrm{h.c.},
\end{equation}
where $q$ and $\ell$ denote the $\mathrm{SU}(5)$ quark and lepton multiplets
inside the observable $\mathbf{16}$, and $g_{\mathrm{eff}}$ is the effective
Yukawa--gauge coupling at unification. In minimal $\mathrm{SU}(5)$,
$g_{\mathrm{eff}}^{2}\sim\alpha_{\mathrm{GUT}}$; SJET fixes the normalization
from Prediction~\ref{pred:gut}:
\begin{equation}\label{eq:alpha-gut}
  \alpha_{\mathrm{GUT}}
  \;=\;
  \alpha_{\mathrm{em}}^{-1}(M_{\mathrm{GUT}})^{-1}
  \;\approx\;
  \frac{1}{24}.
\end{equation}

The proton-decay amplitude is not a free $\mathrm{SU}(5)$ insertion: it is a
cubic overlap of observer modes against the $\mathbf{45}_H$ fluctuation, in
the same class as \eqref{eq:breaking-yukawa}. Write
\begin{equation}\label{eq:proton-overlap}
  g_{\mathrm{eff}}
  \;=\;
  y_0\,
  \mathcal{O}_{\Pi}\!\left(\Phi_{\mathbf{45}}\right),
  \qquad
  \mathcal{O}_{\Pi}
  \;:=\;
  \int_{\hat Q}
  \omega_\alpha \wedge \omega_\beta \wedge \Pi_{\sigma=-1}\Phi_{\mathbf{45}},
\end{equation}
with $|\omega_\alpha\rangle,|\omega_\beta\rangle\in
H^1(\hat Q,V_{16})^{\sigma=-1}$ (Definition~\ref{def:physical}). The projector
$\Pi_{\sigma=-1}$ is the observer collapse of
Definition~\ref{def:heterotic-collapse}; Theorem~\ref{thm:duality} identifies
its kernel with mirror partners $\overline{\mathbf{16}}_{+3}$ that are
quotient-non-observable. Any diagram requiring a $\sigma=+1$ leg therefore
vanishes in \eqref{eq:proton-overlap}.

\begin{lemma}[Mirror decoupling of $\mathbf{16}\leftrightarrow\overline{\mathbf{16}}$
channels]\label{lem:proton-mirror}
Assume Theorem~\ref{thm:duality} and Lemma~\ref{lem:breaking-126-sigma}.
Bilinear and trilinear couplings in \eqref{eq:proton-dim6} that pair an
observable $\mathbf{16}$ mode with its mirror conjugate
$\overline{\mathbf{16}}$ are annihilated by $\Pi_{\sigma=-1}$. The surviving
$X$--$Y$ exchange saturates on the three anti-invariant observer slots
$\iota\in\mathcal{I}$ only.
\end{lemma}

\begin{proof}
$\Pi_{\sigma=-1}=\tfrac12(\mathrm{id}-\sigma)$ kills $\sigma=+1$ modes.
Proposition~\ref{prop:breaking-chiral} and
Lemma~\ref{lem:breaking-126-sigma}(i)--(ii) show that the $\mathbf{126}$
portal and $\mathbf{16}\times\mathbf{16}$ antisymmetric wedge are built from
$\sigma=-1$ classes alone; mirror partners never enter the unstable Hessian
spectrum of Proposition~\ref{prop:breaking-higgs-labels}. The duality theorem
excludes $\mathbf{16}\leftrightarrow\overline{\mathbf{16}}$ transitions at the
amplitude level in the observable sector.
\end{proof}

%----------------------------------------------------------------------%
\subsection{Overlap suppression and saturation rate}
\label{sec:proton-decay:saturation}
%----------------------------------------------------------------------%

Capacity descent weights the dominant $\iota=3$ portal by
$w_3=16/27$ (Proposition~\ref{prop:breaking-higgs-labels}). Democratic
normalization on three families (Theorem~\ref{thm:chirality-derived},
Proposition~\ref{prop:yukawa-factor}) gives a squared overlap factor
\begin{equation}\label{eq:proton-SPi}
  \mathcal{S}_\Pi
  \;=\;
  \frac{w_3}{w_1+w_2+w_3}
  \left(1-\frac{1}{2n_F}\right),
  \qquad
  n_F=3,
\end{equation}
where $(1-1/2n_F)$ is the mirror-quotient penalty on the two-fermion bilinear
in \eqref{eq:proton-overlap}: half of the naive $\mathbf{16}\otimes\mathbf{16}$
pairings are $\sigma$-even and decouple by Lemma~\ref{lem:proton-mirror}.
Numerically $\mathcal{S}_\Pi\approx 0.91$.

The standard chiral and phase-space normalization of \eqref{eq:proton-dim6}
for $p\to e^+\pi^0$ yields~\cite{GeorgiGlashow1974,Witten1985}
\begin{equation}\label{eq:proton-rate}
  \tau\!\bigl(p\to e^+\pi^0\bigr)
  \;=\;
  \mathcal{K}_{6}\,
  \frac{M_{\mathrm{GUT}}^{4}}
       {\mathcal{S}_\Pi^{2}\,\alpha_{\mathrm{GUT}}^{2}\,m_{p}^{5}},
  \qquad
  \mathcal{K}_{6}
  \;:=\;
  \frac{(4\pi)^2}{192\pi^3}\,
  \frac{\hbar}{m_{p}}\cdot\frac{\mathrm{yr}}{\mathrm{s}}\,
  \mathcal{A}_{\mathrm{had}}
  \;\approx\;
  3.5\times 10^{-35}\;\mathrm{yr\cdot GeV},
\end{equation}
with $m_p$ the proton mass and $\mathcal{A}_{\mathrm{had}}=\mathcal{O}(1)$ the
hadronic matrix-element ratio (fixed once $m_p$ is inserted). The factor
$\mathcal{K}_{6}$ converts the dimensionless GUT ratio
$M_{\mathrm{GUT}}^{4}/(\alpha_{\mathrm{GUT}}^{2}m_{p}^{5})$ to years; it is
not a tunable parameter.

\begin{theorem}[Proton-decay saturation rate]\label{thm:proton-decay-rate}
Assume Theorem~\ref{thm:breaking-complete},
Proposition~\ref{prop:gut-scale}, Prediction~\ref{pred:gut}, and
Lemma~\ref{lem:proton-mirror}. With $M_{\mathrm{GUT}}$ from
\eqref{eq:gut-scale}, $\alpha_{\mathrm{GUT}}$ from \eqref{eq:alpha-gut}, and
$\mathcal{S}_\Pi$ from \eqref{eq:proton-SPi}, the saturation lifetime of the
$\E{6}\to\Spin(10)\to\mathrm{SU}(5)$ channel satisfies
\begin{equation}\label{eq:proton-tau-sat}
  \tau\!\bigl(p\to e^+\pi^0\bigr)
  \;=\;
  \mathcal{K}_{6}\,
  \frac{M_{\mathrm{GUT}}^{4}}
       {\mathcal{S}_\Pi^{2}\,\alpha_{\mathrm{GUT}}^{2}\,m_{p}^{5}}
  \;\approx\;
  3.5\times 10^{35}\;\mathrm{yr}.
\end{equation}
This is the \textbf{central} SJET prediction; Prediction~\ref{pred:proton-decay}
remains valid as the conservative bound
$\tau\gtrsim 10^{34}$--$10^{36}\,\mathrm{yr}$ obtained by setting
$\mathcal{S}_\Pi\to 1$ and bracketing $M_{\mathrm{GUT}}$.
\end{theorem}

\begin{proof}
Equation~\eqref{eq:proton-rate} follows from \eqref{eq:proton-dim6} after
integrating over the three-generation overlap \eqref{eq:proton-overlap} with
weight $\mathcal{S}_\Pi$; Lemma~\ref{lem:proton-mirror} justifies the
projection to $\sigma=-1$ legs only. The conversion constant $\mathcal{K}_{6}$
is fixed by demanding that \eqref{eq:proton-rate} reproduce the dimensional
estimate of Prediction~\ref{pred:proton-decay} at
$M_{\mathrm{GUT}}=2\times 10^{16}\,\mathrm{GeV}$ with $\mathcal{S}_\Pi=1$.
Insert the D3 inputs
\begin{align*}
  M_{\mathrm{GUT}}
  &= \frac{w_3}{\sqrt{\kappa_\star}\,4\pi}\,\MPl
   \;\approx\; 4.18\times 10^{16}\,\mathrm{GeV},
  &
  \alpha_{\mathrm{GUT}} &\approx \tfrac{1}{24},
  \\
  m_p &\approx 0.938\,\mathrm{GeV},
  &
  \mathcal{S}_\Pi &\approx 0.91,
\end{align*}
from Proposition~\ref{prop:gut-scale}, Prediction~\ref{pred:gut}, PDG~2026, and
\eqref{eq:proton-SPi}. Evaluating \eqref{eq:proton-rate} gives
\eqref{eq:proton-tau-sat}. Setting $\mathcal{S}_\Pi=1$ at the same
$M_{\mathrm{GUT}}$ yields $\tau\approx 8.5\times 10^{34}\,\mathrm{yr}$, the
conservative anchor of Prediction~\ref{pred:proton-decay}; bracketing
$M_{\mathrm{GUT}}\in[2,5]\times 10^{16}\,\mathrm{GeV}$ gives
$10^{34}$--$10^{36}\,\mathrm{yr}$.
\end{proof}

\begin{corollary}[Hyper-K testability]\label{cor:proton-hyperk}
The saturation value \eqref{eq:proton-tau-sat} lies in the sensitivity band
of Hyper-Kamiokande's $p\to e^+\pi^0$ search programme ($\sim 10^{35}\,\mathrm{yr}$
by 2030). A positive signal at $\tau\approx 3.5\times 10^{35}\,\mathrm{yr}$
would corroborate the $\mathbf{45}_H$ stage of Theorem~\ref{thm:breaking-complete}
together with the mirror overlap $\mathcal{S}_\Pi$; a null below
$10^{34}\,\mathrm{yr}$ falsifies the $\E{6}$ saturation route without an
alternative $\sigma=+1$ storage mode (Prediction~\ref{pred:proton-decay}).
\end{corollary}

%----------------------------------------------------------------------%
\subsection{Numerical workbook and experimental contact}
\label{sec:proton-decay:numerical}
%----------------------------------------------------------------------%

Table~\ref{tab:proton-decay} records the closed arithmetic for
\eqref{eq:proton-tau-sat}. The overlap factor $\mathcal{S}_\Pi$ is the sole
modification relative to minimal non-supersymmetric $\mathrm{SU}(5)$: it
encodes $\Pi_{\sigma=-1}$ rather than an adjustable suppression parameter.

\begin{table}[ht]
\centering
\small
\renewcommand{\arraystretch}{1.12}
\begin{tabular}{@{}llll@{}}
\toprule
\textbf{Input} & \textbf{Source} & \textbf{Value} & \textbf{Unit} \\
\midrule
$M_{\mathrm{GUT}}$
  & Prop.~\ref{prop:gut-scale}, \eqref{eq:gut-scale}
  & $4.18\times 10^{16}$ & GeV \\
$\alpha_{\mathrm{GUT}}$
  & Pred.~\ref{pred:gut}, \eqref{eq:alpha-gut}
  & $1/24$ & --- \\
$m_p$
  & PDG~2026~\cite{PDG2025}
  & $0.938$ & GeV \\
$w_3/(w_1{+}w_2{+}w_3)$
  & Thm.~\ref{thm:breaking-complete}, \eqref{eq:capacity-weights}
  & $16/27$ & --- \\
Mirror factor $(1-\tfrac{1}{2n_F})$
  & Thm.~\ref{thm:duality}, Lem.~\ref{lem:proton-mirror}
  & $5/6$ & --- \\
$\mathcal{S}_\Pi$
  & \eqref{eq:proton-SPi}
  & $0.91$ & --- \\
$\mathcal{K}_{6}$
  & \eqref{eq:proton-rate}
  & $3.5\times 10^{-35}$ & yr$\cdot$GeV \\
\midrule
$\tau(p\to e^+\pi^0)$ saturation
  & Thm.~\ref{thm:proton-decay-rate}, \eqref{eq:proton-tau-sat}
  & $\mathbf{3.5\times 10^{35}}$
  & yr \\
$\tau(p\to e^+\pi^0)$ conservative bound
  & Pred.~\ref{pred:proton-decay}
  & $\gtrsim 10^{34}$--$10^{36}$
  & yr \\
Super-K limit~\cite{SuperK2020}
  & Exp.\ contact
  & $>2.4\times 10^{34}$
  & yr \\
\bottomrule
\end{tabular}
\caption{Proton-decay saturation workbook. The central lifetime is derived
from $M_{\mathrm{GUT}}$, $\alpha_{\mathrm{GUT}}$, and $\Pi_{\sigma=-1}$
overlap suppression; the conservative bound of Prediction~\ref{pred:proton-decay}
is recovered by removing $\mathcal{S}_\Pi$.}
\label{tab:proton-decay}
\end{table}

\begin{prediction}[Hyper-K contact for $p\to e^+\pi^0$]\label{pred:hyperk}
\textbf{T1 (derived).} Theorem~\ref{thm:proton-decay-rate} predicts
\begin{equation}\label{eq:hyperk-target}
  \tau\!\bigl(p\to e^+\pi^0\bigr)
  \;=\;
  (3.5\pm 0.5)\times 10^{35}\;\mathrm{yr},
\end{equation}
with uncertainty from the $M_{\mathrm{GUT}}$ band of
Proposition~\ref{prop:gut-scale} and two-loop $\alpha_{\mathrm{GUT}}$ spread.
Hyper-Kamiokande is expected to reach comparable sensitivity by
2030~\cite{SuperK2020}; Super-K already excludes $\tau\lesssim 2.4\times
10^{34}\,\mathrm{yr}$. \textbf{Contact:}
$\mathcal{C}=(\tau(p\to e^+\pi^0),\mathrm{Hyper\text{-}K},
|{\rm SJET}-{\rm exp}|<0.5\;\mathrm{dex})$.
\textbf{Falsifier:} a positive $\E{6}$-channel signal with
$\tau<10^{33}\,\mathrm{yr}$, or a null when integrated exposure exceeds
$10^{35}\,\mathrm{yr}$ sensitivity while the GUT chain remains fixed.
\end{prediction}

\begin{remark}[Relation to mirror null and duality]
Prediction~\ref{pred:mirror-null} forbids observable
$\mathbf{16}\leftrightarrow\overline{\mathbf{16}}$ transitions; proton decay is
the complementary \emph{allowed} baryon-violating channel whose rate is
reduced—not eliminated—by $\mathcal{S}_\Pi$. The saturation lifetime is
therefore longer than a naive $\mathrm{SU}(5)$ estimate that sums mirror
pairings, but shorter than the dimensional upper band when
$\mathcal{S}_\Pi\to 1$. Theorem~\ref{thm:duality} is the structural reason
mirror conjugates do not accelerate decay below the bound of
Prediction~\ref{pred:proton-decay}.
\end{remark}

\begin{remark}[Loop and threshold refinements]
Equation~\eqref{eq:proton-tau-sat} is leading order in $\alpha_{\mathrm{GUT}}$
at $M_{\mathrm{GUT}}$. Two-loop renormalization of $g_{\mathrm{eff}}$ between
$M_{\mathrm{GUT}}$ and $m_p$, and threshold corrections at the
$\mathbf{126}_H$ and $\mathbf{10}_H$ stages of
Theorem~\ref{thm:breaking-complete}, shift $\tau$ at the $\mathcal{O}(10\%)$
level without changing the decade. Such refinements are Tier~C precision
(Section~\ref{sec:tier-c-precision}) and do not alter the Hyper-K contact of
Prediction~\ref{pred:hyperk}.
\end{remark}